L'espace est muni d'un repère orthonormé \Oijk{}.

On considère :

\begin{itemize}[label=\textbullet~]
	\item les points de l'espace A$(1 ; 0 ; 3)$, B$(2 ; 1 ;-1)$, C$(1 ; 1 ; 1)$ et H$(0 ; 2 ; 1)$;
	\item le vecteur $\vect{n}\begin{pmatrix}4 \\ 4 \\ 2\end{pmatrix}$.
\end{itemize}

\begin{enumerate}
	\item Démontrer que les points A, B et C définissent un plan de l'espace.

	\item Démontrer que le vecteur $\vect{n}$ est un vecteur normal au plan (ABC).

	\item Déterminer une équation cartésienne du plan (ABC).

	\item Vérifier que le point H appartient au plan (ABC).

	\item Déterminer la mesure en degré de l'angle $\widehat{\mathrm{BAH}}$.

	\item Déterminer une représentation paramétrique de la droite $(d)$ orthogonale au plan (ABC) et passant par le point H.

	\item Déterminer les coordonnées du point S$\left(x_{\mathrm{S}} ; y_{\mathrm{S}} ; z_{\mathrm{S}}\right)$ de la droite $(d)$ tel que :
	\begin{itemize}[label=\textbullet~,left=1cm]
		\item la distance entre le point S et le plan (ABC) est 6;
		\item son abscisse $x_{\mathrm{S}}$ est positive.
	\end{itemize}
\end{enumerate}
